\documentclass[11pt]{article}
\usepackage{preamble}
\renewcommand{\answer}[1]{}

\begin{document}

\ladderheading{AP Statistics \textperiodcentered\ Unit 2 \textperiodcentered\ Topic 2.6 \textperiodcentered\ B: 2026-10-26 \textperiodcentered\ E: 2026-11-13}{Which Way Does the Condition Point?}

\focusbanner{TODAY'S PROBLEM}

Suppose a school screens 1{,}000 students for a hypothetical condition. Exactly $5\%$ have the condition. The test has $90\%$ sensitivity: among students with the condition, $90\%$ test positive. Its false-positive rate is $10\%$: among students without the condition, $10\%$ test positive. The implied counts are shown.\par\smallskip \textbf{Mina:} ``The test catches $90\%$ of students who have the condition, so a positive result means a student has a $90\%$ chance of having it.''\par \textbf{Luis:} ``We need to count everyone who tested positive before finding the chance of having the condition.''\par
\begin{center}
\textbf{Screening results (counts)}\par\smallskip
\begin{tabular}{|l|c|c|c|}
\hline
\textbf{Status} & \textbf{Pos.} & \textbf{Neg.} & \textbf{Total} \\ \hline
\textbf{Condition} & 45 & 5 & 50 \\ \hline
\textbf{No condition} & 95 & 855 & 950 \\ \hline
\textbf{Total} & 140 & 860 & 1{,}000 \\ \hline
\end{tabular}
\end{center}
\begin{firsttakebox}{FIRST TAKE --- before the video (5 min)}
Does a positive result mean a $90\%$ chance of having the condition? Choose a claim and cite one table count.
\workspace{0.9in}
\end{firsttakebox}

\begin{callout}{\IconBook\ THE GIVEN EVENT SETS THE DENOMINATOR (keep on the page)}{calloutgreen}
$P(A\mid B)$ means the probability of $A$ \textbf{given} $B$. Restrict attention to the outcomes in $B$:
$P(A\mid B)=\dfrac{P(A\cap B)}{P(B)}=\dfrac{\text{count in both }A\text{ and }B}{\text{count in }B}$.
In general, $P(A\mid B)\neq P(B\mid A)$ because the denominators are different.
\textbf{Sensitivity} is $P(\text{positive}\mid\text{condition})$.
The \textbf{false-positive rate} is $P(\text{positive}\mid\text{no condition})$.
\end{callout}

\newpage
\textbf{Finish after the video:}\par\smallskip

\dokbadge{1}~\textbf{(a)} For $P(\text{condition}\mid\text{positive})$, identify the numerator and denominator counts.\par
\workspace{0.7in}

\dokbadge{2}~\textbf{(b)} Find $P(\text{positive}\mid\text{condition})$ and $P(\text{condition}\mid\text{positive})$. Show each count fraction and percent.\par
\workspace{1.3in}

\dokbadge{3}~\textbf{(c)} Is Mina's claim supported? Use the positive-test counts for students with and without the condition. Explain how reversing given changes the group in the denominator. Write 3--4 sentences.\par
\workspace{2.0in}

\begin{sentenceframebox}\raggedright\textbf{Frame:} The $90\%$ uses the group \blank[1.6in], but a positive result directs us to \blank[1.8in].\end{sentenceframebox}

\begin{sentenceframebox}\raggedright\textbf{Frame:} Among \blank[0.6in] positive results, \blank[0.6in] have the condition, so the chance is \blank[0.8in].\end{sentenceframebox}

\begin{summaryexitbox}{EXIT --- turn this in}
One thing the video changed about my first take: \hrulefill\par
\workspace{0.4in}
\end{summaryexitbox}

\end{document}
